% Chapter III -- Operational Definition of Terms (rev L-104)
% Style: alphabetical; each term defined OPERATIONALLY (how it is measured,
% constructed, or decided in THIS study), not merely conceptually. Every
% definition below is anchored to the concrete artifact that realizes it so
% the panel can trace term -> mechanism -> file.
% REV L-104: cross-references corrected to actual section numbers (shards
% 4.3.5; RAID scoping -> threat T7 in 6.3; headroom -> 5.6; repairs -> 4.6);
% internal repo paths replaced with Appendix pointers.

\chapter{Operational Definition of Terms}

The following terms are defined operationally as they are used throughout
this thesis. Each definition states how the construct was instantiated,
measured, or adjudicated within the executed study, together with the
pipeline component that realizes it, so that terminology remains stable
across the Methodology (Chapter~4), Results (Chapter~5), and the threats-
to-validity analysis (Chapter~6).

\section{Machine-Generated Text (MGT)}
Text produced end-to-end by a large language model conditioned on a prompt
and decoding parameters, without post-generation human rewriting. In this
study the term refers exclusively to generations captured with recorded
(model, prompt-template, decoding-parameter) triples in the surrogate
corpus manifests; any sample whose provenance triple is missing is excluded
from analysis rather than labeled by assumption.

\section{Human-Written Text (HWT)}
Text authored by identified humans under the collection protocols of this
study: (a) participants of the own-corpus arm who provided informed
consent through the consent pipeline, or (b) verified human documents in
the RAID subsample as designated by the benchmark's own labeling. Mixed-
authorship documents---texts begun by humans and completed by machines or
vice versa---are reported as a separate stratum and never collapsed
silently into either class, because the validity audit identifies them as
a confound (threat T3, Chapter~6).

\section{Surrogate Corpus (surrogate\_v2)}
The synthetic training corpus constructed for this study: machine-generated
documents sampled across a controlled grid of generator families, prompt
templates, topics, and lengths, assembled by the surrogate\_v2 pipeline and
governed by its licensing gate. ``Surrogate'' signals the design decision
that detector training uses proxy generators rather than the target
generators evaluated at test time, making the surrogate--real gap a first-
class object of measurement (Section~4.3.1; conclusions in Section~6.2).

\section{Own Corpus}
The corpus of participant-authored texts collected directly for this study
through intake forms, consent capture, and automated guard checks
(personally identifiable information scrubbing, scope validation) prior to
admission. Admission of any batch into the experimental splits required the
feasibility gate (below) to pass; the frozen state of that gate defines the
exact data snapshot analyzed in Chapter~5.

\section{Own-Corpus Feasibility Gate}
A predeclared go/no-go rule evaluating whether the accumulated own corpus
meets minimum adequacy criteria (volume, class balance, metadata
completeness, guard-pass rate) before any model is trained on it. The gate
was frozen at the date recorded in its specification memo; ``frozen''
means the criteria and threshold values were fixed before outcome
inspection to prevent post hoc data-dependent relaxation. Its outcome and
the freeze date constitute part of the study's predeclared-analysis record
(freeze memo reproduced in Appendix~E). This record is analogous to, but
weaker than, formal preregistration; it is not claimed as the latter.

\section{RAID Subsample}
The stratified subset of the RAID benchmark used for external, out-of-
training-distribution validation. Subsampling was necessitated by compute
budget and implemented as stratification over RAID's native domains,
decoding strategies (greedy vs.\ sampled), and attack conditions, with the
stratum manifest retained for reproducibility (Appendix~D). All claims
about generalization in this thesis are scoped to this subsample (threat
T7, Section~6.3).

\section{Recovered Device Shards (phd\_rescued)}
Document shards recovered from participant devices during a salvage
operation, admitted into the study only after passing a provenance review
documented in a chain-of-custody record (Appendix~F). Because their
collection pathway differs from the prospective consent pipeline, they are
treated as a distinct quarantined data tier with its own robustness checks
(Section~4.3.5), and headline results are reported both with and without
them (selection threat T5, Section~6.3).

\section{Detector}
Any scoring function $s(x)\!\in\![0,1]$ mapping a document to the
probability that it is machine-generated. The study distinguishes:
(a) the \emph{surrogate-trained classifier}, the principal system under
evaluation; (b) \emph{zero-shot statistical baselines} requiring no
training data; and (c) \emph{embedding-based baselines} using frozen
representations. Detector identity is always reported jointly with its
training data tier.

\section{Detection Performance (AUROC / AUPRC)}
Area under the receiver operating characteristic curve and area under the
precision--recall curve computed on held-out splits, reported with bootstrap
confidence intervals. AUROC is the primary headline metric; AUPRC
accompanies it because class prevalence varies across strata and AUROC
alone can overstate operating utility.

\section{Reliability / Calibration}
The agreement between predicted probabilities and empirical error rates,
assessed via reliability diagrams and Expected Calibration Error (ECE) at
fixed binning. A detector is called \emph{well-calibrated on a stratum}
when its reliability curve lies within the predeclared tolerance band of
the diagonal on that stratum [[FILL: tolerance value from calibration
configuration; do not state from memory]]. Calibration is treated as a
co-equal outcome, since deployment-relevant thresholds presuppose
trustworthy scores.

\section{Paired Delta}
The difference in a metric between two systems measured on identical
evaluation items (e.g., surrogate-trained classifier minus zero-shot
baseline on the same RAID stratum), summarized with the paired bootstrap
distribution over items. Pairing on items, rather than comparing
independent aggregates, is the inferential backbone of every ``X improves
over Y'' claim in Chapter~5.

\section{Headroom}
For a given length bucket or condition, the gap between achieved detection
performance and an empirical upper-reference derived from the strongest
observed system on that slice. This reference is a measured bound, not an
oracle ceiling; headroom therefore quantifies remaining distance to the
best demonstrated performance, not to perfect detection. The headroom-by-
length audit decomposes overall error into slices, revealing where
apparent aggregate competence conceals localized failure (Section~5.6);
``residual headroom'' denotes gap remaining after the repair attempts
described in Section~4.6.

\section{Evasion (Paraphrastic Attack)}
The transformation of MGT by paraphrasing models or prompts intended to
defeat detection while preserving meaning, as instantiated by the evader
suite carried over from the idea-phase work and by the attack conditions
present natively in RAID. Robustness claims always specify attacker
strength and whether the detector saw attacked text in training.

\section{Linguistic Splitting}
The idea-phase construct underlying feature design: partitioning documents
by measurable linguistic dimensions (register, discourse structure,
lexical diversity bands) so that detector behavior can be profiled per
split rather than aggregated. It survives in this thesis as the slicing
scheme used in the headroom and breakdown audits.

\section{Intrinsic Dimension (of Representations)}
The estimated dimensionality of the embedding manifold occupied by a
document set, computed with a fixed estimator in the project's metrics
toolchain (settings frozen in the run configurations) as a descriptive
statistic contrasting HWT and MGT manifolds. It is used interpretively
(Chapter~5), not as a standalone classifier.

\section{Consent Pipeline / Participant-Data Guard}
The paired mechanisms governing ethical data handling: consent capture
(versioned forms, timestamped receipts) and automated screening that
blocks ingestion of material failing privacy or scope checks. Their joint
pass is a precondition for any text entering the own-corpus tier.

\section{Composition Provenance (Manuscript-Level)}
The per-section record of how thesis prose was produced (typed by the
researcher, manually revised, or machine-assisted), shipped as
Appendix~H. This term exists to draw a bright line between the thesis's
research objects (detectors, evaders) and its own writing process:
evasion-oriented tooling is applied to experiment data only, never to
manuscript text.